% Image credits for the photographs and AI illustrations of Book 1 (Primary and Middle School).
% Machine-readable license records live in images/book1/CREDITS.md; keep the
% two files in sync. The Book 1 writer owns this file: add one \item per
% photograph (inside the itemize, by chapter) as photographs land.
\chapter*{Verantwoording van de afbeeldingen}

De tekeningen, kaarten en grafieken in dit boek zijn origineel; kaarten en
grafieken die op gepubliceerde gegevens berusten, noemen hun bron in het
onderschrift of in deze lijst. De foto's worden gebruikt onder hun
respectieve vrije licenties, met dank aan hun makers; de volledige
bronlinks en licentie-URL's van elke foto staan in
\texttt{images/book1/CREDITS.md} in de repository van het boek.

\bigskip
Naast de foto's en de originele tekeningen zijn sommige figuren door AI
gegenereerde illustraties, gemaakt met de beeldgeneratie van OpenAI op
basis van prompts die de redactie van het boek heeft geschreven, en vóór
opname op wetenschappelijke juistheid gecontroleerd. Ze tonen alleen
algemene scènes, nooit een echte plek, een echte persoon of een echt
exemplaar. De volledige prompt van elke gegenereerde afbeelding staat in
\texttt{images/book1/ai/PROMPTS.md} in de repository.

\bigskip
\noindent Foto's en gegevens, per hoofdstuk:
\begin{itemize}\small
  \item Hfst.~\ref{ch:g2:water-on-earth} (water op Aarde): de Aarde gezien vanuit Apollo~17, 7~december 1972 (``The Blue Marble''); NASA, publiek domein (via Wikimedia Commons). De 100 vakjes zijn berekend uit het reliëf van NOAA ETOPO1; de 100 emmers uit de tabel van de Water Science School van de U.S. Geological Survey (naar Shiklomanov, 1993).
  \item Hfst.~\ref{ch:g3:stones-pebbles-sand} (stenen, kiezels en zand): een droge stapelmuur van graniet; foto van Tony Atkin, Geograph, CC BY-SA 2.0 (via Wikimedia Commons). Daken van stenen platen in Zwitserland; foto van Michael C. Rygel, CC BY-SA 3.0 (via Wikimedia Commons).
  \item Hfst.~\ref{ch:g3:continents-and-oceans} (de globe): de twee globes en de wereldkaart zijn getekend met de kustlijnen van Natural Earth (publiek domein).
  \item Hfst.~\ref{ch:g4:clouds-rain-snow} (wolken, regen en sneeuw): cumulus boven Kansas (S.~Corfidi), plukjes cirrus (R.~F.~Kresge), een laag stratus rond bergtoppen, een onweer boven de woestijn van Arizona (collectie van B.~Hooke), een zeer grote hagelsteen; alle NOAA Photo Library, publiek domein (via Wikimedia Commons).
  \item Hfst.~\ref{ch:g4:volcanoes-earthquakes-first-look} (vulkanen en aardbevingen): lava die bij de Kilauea de oceaan in stroomt, USGS Hawaiian Volcano Observatory; Mount St.~Helens, 18~mei 1980, foto van Austin Post, USGS; een huis dat door de aardbeving van Loma Prieta in 1989 van zijn fundering is geworpen, foto van J.~K.~Nakata, USGS; alle publiek domein (via Wikimedia Commons). De kaart gebruikt de lijst van holocene vulkanen van het Global Volcanism Program van het Smithsonian en de aardbevingscatalogus van de U.S. Geological Survey.
  \item Hfst.~\ref{ch:g4:world-physical-map} (de natuurkundige kaart): de Himalaya en de Amazone, gefotografeerd vanuit het Internationaal Ruimtestation, NASA, publiek domein (via Wikimedia Commons). De natuurkundige kaart is getekend met het reliëf van NOAA ETOPO1, met rivieren, meren en kustlijnen van Natural Earth (publiek domein).
  \item Hfst.~\ref{ch:g5:water-shapes-the-land} (water vormt het land): de rivier de Mamoré, gefotografeerd vanuit het Internationaal Ruimtestation (ISS007-E-10797), en de Nijldelta, gezien door het MODIS-instrument van de satelliet Terra, NASA, publiek domein (via Wikimedia Commons).
  \item Hfst.~\ref{ch:g5:weather-and-climate} (weer en klimaat): de Sonorawoestijn, U.S. National Park Service, Saguaro National Park, publiek domein (via Wikimedia Commons). De klimaatdiagrammen zijn getekend met de U.S. Climate Normals 1991--2020 van NOAA NCEI (publiek domein).
  \item Hfst.~\ref{ch:g5:natural-hazards} (natuurgevaren): het oog van orkaan Isabel, gefotografeerd vanuit het Internationaal Ruimtestation door Ed Lu, NASA, publiek domein; de aardverschuiving van 2014 in de staat Washington, Mark Reid, U.S. Geological Survey, publiek domein (beide via Wikimedia Commons).
  \item Hfst.~\ref{ch:g6:earth-a-planet} (de Aarde, een planeet): Venus vanuit Mariner~10 (NASA/JPL-Caltech); de Aarde vanuit Apollo~17 (NASA); \emph{Earthrise} vanuit Apollo~8 (NASA, Bill Anders); Mars en Jupiter gezien door de Hubble-ruimtetelescoop (NASA, ESA, het Hubble Heritage Team en A.~Simon); alle publiek domein (via Wikimedia Commons). De \emph{Pale Blue Dot}, gezien door Voyager~1 in 1990, versie van 2020 (NASA/JPL-Caltech, PIA23645, publiek domein). Gegevens over de planeten uit de planetary fact sheet van NASA.
  \item Hfst.~\ref{ch:g6:latitude-longitude} (breedte en lengte): de H4 van Harrison in het National Maritime Museum, Greenwich, foto van Colonel Warden, CC BY-SA 3.0 (via Wikimedia Commons). Kustlijnen en steden van Natural Earth (publiek domein).
  \item Hfst.~\ref{ch:g6:rocks-and-minerals} (gesteenten en mineralen): graniet, basalt, fossielhoudende kalksteen, zandsteen, marmer, schalie, een kwartskristal en een plaatje muscoviet (mica), foto's van James St.~John, CC BY 2.0; een leien dak, foto van Øyvind Holmstad, CC0 (alle via Wikimedia Commons).
  \item Hfst.~\ref{ch:g6:inside-the-earth} (het binnenste van de Aarde): Inge Lehmann in 1932, foto van Even Neuhaus, Koninklijke Bibliotheek van Denemarken, publiek domein (via Wikimedia Commons). De lagen, dichtheden en golfsnelheden zijn ontleend aan het Preliminary Reference Earth Model (A.~M. Dziewonski en D.~L. Anderson, 1981).
  \item Hfst.~\ref{ch:g6:water-underground} (water onder de grond): Big Spring, Missouri, foto van Johnston9494, CC BY-SA 4.0 (via Wikimedia Commons). Porositeiten naar R.~C. Heath, \emph{Basic Ground-Water Hydrology}, U.S. Geological Survey (1983); aandelen zoet water van de Water Science School van de USGS.
  \item Hfst.~\ref{ch:g7:plate-tectonics} (platentektoniek): de zuidelijke continenten weer samengevoegd, met de fossielenbanden van \emph{Cynognathus}, \emph{Lystrosaurus}, \emph{Mesosaurus} en \emph{Glossopteris}; U.S. Geological Survey, \emph{This Dynamic Earth}, publiek domein.
  \item Hfst.~\ref{ch:g7:plate-tectonics}: de San-Andreasbreuk door de Carrizo Plain; foto van Robert E. Wallace, U.S. Geological Survey, publiek domein.
  \item Hfst.~\ref{ch:g7:plate-tectonics}: plaatgrenzen naar P.~Bird (2003), GeoJSON van H.~Ahlenius (ODC-By 1.0); kustlijnen van Natural Earth (publiek domein); aardbevingen uit de catalogus van de U.S. Geological Survey; vulkanen van het Global Volcanism Program van het Smithsonian.
  \item Hfst.~\ref{ch:g7:volcanoes} (vulkanen): Mount St.~Helens vanaf Johnston Ridge op 17~mei en 10~september 1980, foto's van Harry Glicken, U.S. Geological Survey; een lavafontein op de Kilauea, 1983, en de Mauna Loa vanuit de lucht, 1985, foto's van J.~D. Griggs, U.S. Geological Survey; de eruptiekolom van de Pinatubo, 12~juni 1991, foto van Dave Harlow, U.S. Geological Survey; alle publiek domein (via Wikimedia Commons). De kaart en de tellingen gebruiken de lijst van holocene vulkanen van het Global Volcanism Program van het Smithsonian; plaatgrenzen naar P.~Bird (2003).
  \item Hfst.~\ref{ch:g7:earthquakes} (aardbevingen): een dorp bij de kust van Sumatra na de tsunami van 26~december 2004, foto van de U.S. Navy door Philip A. McDaniel; de breukranden van de aardbeving van Borah Peak in 1983, foto van R.~E. Wallace, U.S. Geological Survey; de weg die in het Hebgenmeer is weggezakt, 1959, foto van I.~J. Witkind, U.S. Geological Survey; alle publiek domein (via Wikimedia Commons). De tellingen en diepten gebruiken de aardbevingscatalogus van de U.S. Geological Survey; de intensiteitstabel is ingekort naar de gewijzigde schaal van Mercalli van de U.S. Geological Survey.
  \item Hfst.~\ref{ch:g7:sediments-and-fossils} (sedimenten en fossielen): de Grand Canyon vanaf Hermits Rest, foto van de NPS door Kristen M. Caldon, publiek domein; een ammoniet, een trilobiet en ribbels in Capitol Reef National Park, foto's van James St.~John, CC BY 2.0; dinosaurussporen bij Dinosaur Ridge, foto van Dexter Perkins, CC0 (alle via Wikimedia Commons).
  \item Hfst.~\ref{ch:g7:geological-time} (de ouderdom van de Aarde): James Hutton, portret door Henry Raeburn, publiek domein (via Wikimedia Commons). De ouderdommen en kleuren van de tijdschaal zijn die van de International Chronostratigraphic Chart (International Commission on Stratigraphy).
  \item Hfst.~\ref{ch:g8:the-atmosphere} (de atmosfeer): het oplaten van een weerballon in Topeka, Kansas, foto van Stephen Corfidi, NOAA/National Weather Service; de rand van de Aarde bij zonsondergang vanuit het Internationaal Ruimtestation, NASA, ISS Expedition~23; beide publiek domein (via Wikimedia Commons). De profielen zijn berekend uit de \emph{U.S. Standard Atmosphere, 1976} (NOAA, NASA en U.S. Air Force).
  \item Hfst.~\ref{ch:g8:winds-global-circulation} (wind en de algemene circulatie): Noord- en Zuid-Amerika gezien door de satelliet GOES-16 op 26~juli 2022, NOAA, publiek domein (via Wikimedia Commons). De kustlijnen van de kaart van de passaatwinden zijn van Natural Earth (publiek domein).
  \item Hfst.~\ref{ch:g8:fronts-and-weather-maps} (fronten en weerkaarten): een depressie bij IJsland, gezien door de satelliet Terra, NASA/GSFC, MODIS Rapid Response Team, Jacques Descloitres, publiek domein (via Wikimedia Commons). Kustlijnen van de luchtmassakaart van Natural Earth (publiek domein).
  \item Hfst.~\ref{ch:g8:the-oceans} (de oceanen): HMS \emph{Challenger} maakt zich klaar voor een dieptepeiling, 1872, gravure uit de verslagen van de Challenger-expeditie, overgenomen in W.~A. Herdman, \emph{Founders of Oceanography} (1923); de Golfstroom, gezien door het MODIS-instrument van de satelliet Aqua, NASA Earth Observatory; beide publiek domein (via Wikimedia Commons). De profielen van de oceaanbodem zijn getekend met het reliëf van NOAA ETOPO1; kustlijnen van de kaarten van de gyres en de boeien van Natural Earth (publiek domein).
  \item Hfst.~\ref{ch:g8:ice-on-earth} (ijs op Aarde): de Pedersengletsjer in 1990 en 2016, foto's van M.~Tetreau en D.~Kurtz, U.S. National Park Service; de rand van een ijsplaat bij het Antarctisch Schiereiland, NASA, Operation IceBridge (M.-J. Viñas); alle publiek domein (via Wikimedia Commons). De grafiek van het zee-ijs gebruikt de Sea Ice Index van het NSIDC, versie~4 (F.~Fetterer en anderen, 2025, doi:10.7265/a98x-0f50).
  \item Hfst.~\ref{ch:g9:climate-past-and-present} (het klimaat van de Aarde, vroeger en nu): een stuk van de GISP2-ijskern uit Groenland, NOAA Paleoclimatology, publiek domein (via Wikimedia Commons). De reeksen zijn getekend met NASA GISTEMP versie~4, de kooldioxidereeks van Mauna Loa (NOAA Global Monitoring Laboratory en Scripps Institution of Oceanography), de temperatuur van EPICA Dome~C (J.~Jouzel en anderen, 2007), de samengestelde kooldioxidereeks uit Antarctische ijskernen (B.~Bereiter en anderen, 2015) en de zeespiegel van het NOAA Laboratory for Satellite Altimetry; de uitspraken over waarschijnlijkheid zijn die van het zesde beoordelingsrapport van het IPCC, Werkgroep~I, Samenvatting voor beleidsmakers (2021).
  \item Hfst.~\ref{ch:g9:natural-risks} (leven met natuurrisico's): de gevarenkaart van de Nevado del Ruiz, uit T.~L. Wright en T.~C. Pierson, \emph{Living with Volcanoes}, U.S. Geological Survey Circular 1073 (1992); een tsunamiwaarschuwingsboei, geïnspecteerd door de bemanning van de USCGC \emph{William Hart}, U.S. Coast Guard, 2021; beide publiek domein (via Wikimedia Commons).
  \item Hfst.~\ref{ch:g9:earth-resources} (grondstoffen uit de Aarde): de mijn van Chuquicamata, gezien door ASTER, NASA\slash METI\slash AIST\slash Japan Space Systems en het U.S./Japan ASTER Science Team, 2007, publiek domein; een steenkoollaag in een weginsnijding in Belmont County, Ohio, foto van James St.~John, CC BY 2.0 (beide via Wikimedia Commons). De concentratiefactoren komen uit D.~Perkins en anderen, \emph{Mineralogy}, en S.~Earle, \emph{Environmental Geology}.
  \item Hfst.~\ref{ch:g9:landforms} (landvormen): de Seven Sisters, foto van Arild Vågen, 2017, CC BY-SA 4.0; sikkelduinen aan de kust van Zuid-Brazilië, foto van een NASA-astronaut vanuit het Internationaal Ruimtestation, 2017, publiek domein; stalactieten in Carlsbad Cavern, foto van de National Park Service door Daniel Leifheit, 2018, publiek domein (alle via Wikimedia Commons).
  \item Hfst.~\ref{ch:g9:maps-of-the-earth} (kaarten van de hele Aarde): de delta van de Lena, Landsat-beeld van 2000, USGS EROS Data Center via NASA Visible Earth, publiek domein; een GPS-satelliet, voorstelling van een NASA-kunstenaar, publiek domein (beide via Wikimedia Commons). Kustlijnen op de wereldkaarten van Natural Earth (publiek domein).
\end{itemize}
